\chapter{Literature Review}


Academic interest in the evaluation of rootkit detection rises once newer technologies are introduced \& techniques are discovered - such as with eBPF \& io\_uring.
Consequently in the last 10 years there is a sizeable amount of documentation on the various methodologies utilised by those wishing to evaluate a rootkit detector.

This literature review will attempt to cover \& summise the existing methodologies in the literature for evaluating rootkit detector effectiveness.
Additionally, It will cover the results found in these studies as useful datapoints for comparison to the results found in this project.

Finally, it will also discuss any relevant literature regarding the development of rootkits, aiming to give a deeper technical explanation of their ability to evade detection.

\section{Existing Methodologies For Evaluating Rootkits}

\section{Results of Current Evaluations}

\section{Modern Rootkit Design \& Detector Engineering}
